% This LaTeX document needs to be compiled with XeLaTeX.
\documentclass[10pt]{article}
\usepackage[utf8]{inputenc}
\usepackage{amsmath}
\usepackage{amsfonts}
\usepackage{amssymb}
\usepackage[version=4]{mhchem}
\usepackage{stmaryrd}
\usepackage{bbold}
\usepackage[fallback]{xeCJK}
\usepackage{polyglossia}
\usepackage{fontspec}





\usepackage{mdframed}
\usepackage{xcolor}
\definecolor{redcolor}{rgb}{1.0, 0.0, 0.0}            % {\color{redcolor}...}
\definecolor{greencolor}{rgb}{0.3, 1.0, 0.3}          % {\color{greencolor}...}
\definecolor{bluecolor}{rgb}{0.0, 0.0, 1.0}           % {\color{bluecolor}...}
\definecolor{yellowcolor}{rgb}{1.0, 1.0, 0.0}         % {\color{yellowcolor}...}
\definecolor{cyancolor}{rgb}{0.0, 1.0, 1.0}           % {\color{cyancolor}...}
\definecolor{magentacolor}{rgb}{1.0, 0.0, 1.0}        % {\color{magentacolor}...}
\definecolor{blackcolor}{rgb}{0, 0, 0}                % {\color{blackcolor}...}
\definecolor{graycolor}{rgb}{0.2, 0.2, 0.2}           % {\color{graycolor}...}
\definecolor{darkgreen}{RGB}{0,128,0}                 % {\color{darkgreen}...}
\usepackage{ifthen}


% ------------------------------------------------------------------------------
% \begin{define}{<title>}[<label>][<color>] … \end{define}
% ------------------------------------------------------------------------------
% DESCRIPTION:
%   This environment creates a framed definition or theorem with colored background
%   and border. It provides a visually distinct way to present important concepts.
%
% FEATURES:
%   - Colored frame and background (10% opacity)
%   - Section symbol (§) followed by title
%   - Configurable colors for customization
%   - Proper spacing and typography
%   - Optional cross-reference label support
%
% PARAMETERS:
%   {<title>}    Title of the definition/theorem (required)
%   [<label>]    Label for cross-referencing (optional, e.g. "def:continuity").
%                If provided, enables \nameref{<label>} and \pageref{<label>}.
%   [<color>]    Background and border color (optional, defaults to graycolor)
%
% DEPENDENCIES:
%   \usepackage{mdframed}
%   \usepackage{xcolor}
%   \usepackage{ifthen}
%
% USAGE EXAMPLES:
%   % Case 1: Basic usage - only title, default gray color, no cross-reference
%   \begin{define}{Definition of Continuity}
%     A function $f$ is continuous at a point $a$ if for every $\varepsilon > 0$, 
%     there exists a $\delta > 0$ such that $|f(x) - f(a)| < \varepsilon$ whenever $|x - a| < \delta$.
%   \end{define}
%
%   % Case 2: With cross-reference label but default color
%   \begin{define}{Fundamental Theorem of Calculus}[thm:fundamental]
%     If $f$ is continuous on $[a,b]$ and $F$ is an antiderivative of $f$, then
%     $\int_a^b f(x) \, dx = F(b) - F(a)$.
%   \end{define}
%
%   % Case 3: With both cross-reference label and custom color
%   \begin{define}{Euler's Identity}[eq:euler][blue]
%     $e^{i\pi} + 1 = 0$, where $e$ is Euler's number, $i$ is the imaginary unit, 
%     and $\pi$ is the ratio of a circle's circumference to its diameter.
%   \end{define}
%
%   % Case 4: With custom color but no cross-reference
%   \begin{define}{Pythagorean Theorem}[][green]
%     In a right triangle, the square of the hypotenuse equals the sum 
%     of squares of the other two sides: $a^2 + b^2 = c^2$.
%   \end{define}
%
%   % Case 5: Referencing the labeled definitions
%   As we saw in \nameref{thm:fundamental}, integration and differentiation 
%   are inverse operations. The beautiful \nameref{eq:euler} connects 
%   five fundamental mathematical constants.
% ------------------------------------------------------------------------------
\makeatletter % "ENABLE" internal \@... macros
\NewDocumentEnvironment{define}{ m O{} O{graycolor} }
  {
  \vspace{2pt}
  \begin{mdframed}[backgroundcolor=#3!10, linecolor=#3]
  \color{#3}
  % Create cross-reference if label is provided
  \ifthenelse{\equal{#2}{}}{}{%
    \phantomsection
    \def\@currentlabelname{#1}%
    \label{#2}%
  }%
  \textbf{\S~#1}\hspace{6pt}  % Start symbol for definition/theorem with a title
  \vspace{4pt}
  }
{
  \end{mdframed}
}
\makeatother % "DISABLE" internal \@... macros

%%%%%%%%%%%%%%%%%%%%%%%%%%%%%%%%%%%%%%%
% Prefix	Indexed Object Type	    
%%%%%%%%%%%%%%%%%%%%%%%%%%%%%%%%%%%%%%%
% cha:    Chapter (\chapter)
% sec:    Section (\section)
% subsec: Subsection (\subsection)
% prob:	  Problem/Exercise (\item)
% fig:    Figure (figure)
% tab:    Table (table)
% eq:     Equation (equation)
% app:	  Appendix (\appendix)
% algo:   Algorithm (algorithm)
% def:	  Definition (definition)
% prop:   Proposition (proposition)
% thm:    Theorem (theorem)
% lem:	  Lemma (lemma)
%%%%%%%%%%%%%%%%%%%%%%%%%%%%%%%%%%%%%%%









\begin{document}

L8: Corollaries to Fermat's Little Theorem and Psendoprime numbers\\
Corollary 2.14: Let $p$ be a prime and let $a \in \mathbb{Z}$. If $p+a$, then $a^{p-2}$ is the inverse of $a$ modulo $p$.\\
Proof: By Fermat's Little Theorem, we have that $a^{p-1} \equiv 1(\bmod p)$. Consequently, $a a^{p-2} \equiv(\bmod p)$, and so $a^{p-}$\\
as desired. as desired.\\
Remark: Corollary 2.14 is useful in solving linear congruences in one variable of the form $a x \equiv b(\bmod p)$ where $p$ is a prime and $p+a$ i.e. multiply both sides by $a^{p-2}$ to solve. Corollary 2.15: Let $p$ be a prime and let $a \in \mathbb{Z}$. Then $a^{p} \equiv a(\bmod p)$.\\
Proof: If pta, then the Corollary follows from Fermat's Little Theorem on multiplication of both sides of $a^{p-1} \equiv 1(\bmod p)$ by $a$.

Corollary 2.16: Let $p$ be a prime. Then $2^{P} \equiv 2(\bmod p)$.\\
Proof: Put $a=2$ in the statement of Corollary 2.15.

\begin{itemize}
  \item Is the converse of Corollary 2.16 true?\\
i.e. If $n \in \mathbb{Z}, n>1$ and $2^{n} \equiv 2(\bmod n)$, then is $n$ necessarily a prime? If so, this would give a necessary and sufficient condition for an integer $n>1$ to be prime (similar to Wilson's Theorem and its converse). Chinese mathematicians more than 2000 years ago claimed such a characterization. Unfortunately, the following counter example disproves their claim.\\
Example: Consider $n=341$. Observe that $n=341 =11.31$, so not a prime. To prove that $2^{341} \equiv 2(\bmod 341)$, it suffices to prove that $2^{341} \equiv 2(\bmod 11)$ and $2^{341} \equiv 2(\bmod 31)$, since\\
$(11,31)=1$. We have that
\end{itemize}

$$
\begin{aligned}
2^{341} & \equiv\left(2^{10}\right)^{34} \cdot 2(\bmod 11) \\
& \equiv 1^{34} \cdot 2(\bmod 11) \text { by Fermat Little Thm } \\
& \equiv 2(\bmod 11),
\end{aligned}
$$

and

$$
\begin{aligned}
2^{341} & \equiv\left(2^{30}\right)^{11} \cdot 2^{11}(\bmod 31) \\
& \equiv 1^{11} \cdot 2^{11}(\bmod 31) \text { by Fermat Little Thm } \\
& \equiv\left(2^{5}\right)^{2} \cdot 2(\bmod 31) \\
& \equiv 1^{2} \cdot 2 \\
& \equiv 2(\bmod 31)
\end{aligned}
$$

Thus $2^{341} \equiv 2(\bmod 341)$, but 341 is not prime.

Definition: Let $n$ be a composite integer. If $2^{n} \equiv 2(\bmod n)$, then $n$ is said to be a psendoprime.

\begin{itemize}
  \item 341 is the first psendoprime, followed by 561 and 645. There are both infinitely many\\
even and odd pseudoprime numbers.
\end{itemize}

\section*{Euler's Theorem}
\begin{itemize}
  \item Enler's Theorem generalizes Fermat's Little Theorem: if $m$ is any positive integer and $a$ is an integer such that $(a, m)=1$, then Euler's Theorem tells you a particular power of a that is guaranteed to be congruent to 1 modulo m.
\end{itemize}

Definition: Let $n \in \mathbb{Z}$ with $n>0$. The Euler $\phi$ - function, denoted $\phi(n)$, is defined by $\phi(n):=\mid\{x \in \mathbb{Z}: \mid \leqslant x \leq n$ and $(x, n)=1\} \mid$.\\
Thus $\phi(n)$ counts the number of positive integers less than or equal to $n$ that are relatively prime to $n$.

$$
\begin{array}{rlrl}
\text { Example: } & \cdot & \phi(12)=4 & \\
\cdot & \phi(14)=6 & & (1,5,7,11) \\
\cdot & \phi(p)=p-1 & \text { for } p \text { a prime. } \\
& & (1,2, \ldots, p-1) .
\end{array}
$$

\section*{Theorem 2.17 (Euler's Theorem)}
Let $a, m \in \mathbb{Z}$ with $m>0$. If $(a, m)=1$, then

$$
a^{\phi(m)} \equiv 1(\bmod m) .
$$

Remark: Fermat's Little Theorem is a special case of Euler's Theorem i.e. take $m=p$, a prime number, and then we have $\phi(p)=p-1$. The Statement then becomes $a^{p-1} \equiv 1(\bmod p)$.\\
Proof (of Euler's Theorem): Let $r_{1}, r_{2}, \ldots, r_{\phi(m)}$ be the $\phi(m)$ distinct positive integers not exceeding $m$ such that $\left(r_{i}, m\right)=1, i=1,2, \ldots, \phi(m)$. Consider the $\phi(m)$ integers given by $r_{1} a, r_{2} a, \ldots$, $r_{\phi(m)} a$. Note that $\left(r_{i} a, m\right)=1$ for $i=1,2, \ldots, \phi(m)$ Otherwise, if $(r i a, m)>1$ for some $i$, then there is a prime divisor $p$ of ( $r_{i} a, m$ ) by Lemma 1.5 from which $p \mid r i a$ and $p \mid m$. Now, $p$ |ria implies that $p$ |ri or $p l a$ by Lemma 1.14; so either we have $P / r_{i}$ and p/m or p/a and p/m. Both of these\\
cases are impossible since $(r i, m)=1$\\
and $(a, m)=1$ respectively. Also note that no two of the $r_{1}, r_{2}, \ldots, r_{\phi(m)} a$ are congruent modulo $m$. Since $(a, m)=1$, the inverse of a modulo $m$, say $a^{\prime}$, exists by Corollary 2.8. Hence if $r_{i} a \equiv r_{j} a(\bmod m)$ with $i \neq j$, then $r_{i} a q^{\prime} \equiv r_{j} a q^{\prime}(\bmod m)$, and So $r_{i} \equiv r_{j}(\bmod m)$, which is impossible. Thus the least non-negative vesidues modulo $m$ of $r_{1}, a, r_{2} a, \ldots, r_{\phi(m)} a$, taken in some order, must be $r_{1}, r_{2}, \cdots, r_{\phi(m)}$. Then\\
$\left(r_{1} a\right)\left(r_{2} a\right) \cdots\left(r_{\phi(m)} a\right) \equiv r_{1} r_{2} \cdots r_{\phi \mid m}(\bmod m)$,\\
or equivalently,

$$
a^{\phi(m)} r_{1} r_{2} \cdots r_{\phi(m)} \equiv r_{1} r_{2} \cdots r_{\phi(m)}(\bmod m)
$$

Now $m \mid a^{\phi(m)} r_{1} r_{2} \cdots r_{\phi(m)}-r_{1} r_{2} \cdots r_{\phi(m)}$ implies that $m \mid r_{1} r_{2} \cdots r_{\phi(m)}\left(a^{\phi(m)}-1\right)$. Since $\left(r_{1} r_{2} \cdots r_{\phi(m)}, m\right)=1$, we have $m \mid a^{\phi(m)}-1$, and $a^{\phi(m)} \equiv 1(\bmod m)$, as desired.

Definition: Let $m$ be a positive integer. A set of $\phi(m)$ integers such that each element of the set is relatively prime to $m$, and no two elements of the set are congruent modulo $m$ is said to be a reduced residue system modulo $m$.\\
Example: $\{1,5,7,11\}$ is a reduced residue system modulo 12

\begin{itemize}
  \item $\{1,3,5,9,11,13\}$ is a reduced vesidue system modulo 14
  \item $\{1,2, \ldots, p-1\}$ is a reduced residne System modulo $p$.\\
Idea 2.18: Let $m$ be a positive integer and let $\left\{v_{1}, r_{2}, \ldots, r_{\phi(m)}\right\}$ be a reduced vesidue system modulo $m$. If $a$ is an integer with $(a, m)=1$, then the set $\left\{a r_{1}, a r_{2}, \ldots, a r_{\phi(m)}\right\}$ is a reduced residue system modulo $m$.
  \item Idea 2.18 was the key to proving Euler's Theorem.
\end{itemize}

Example: $\{1,5,7,11\}$ reduced vesidue

\section*{(8)}
system modulo 12. Since $(5,12)=1$, the set $\{1.5,5.5,7.5,11.5\}=\{5,25,35,55\}$ is also a reduced residue system modulo 12.

Corollary 2.19 (to Euler's Theorem): Let $a, m \in \mathbb{Z}$ with $m>0$. If $(a, m)=1$, then $a^{\phi(m)-1}$ is the inverse of a modulo $m$.\\
Proof: By Euler's Theorem, we have that $a^{\phi(m)} \equiv 1(\bmod m)$. Hence $a a^{\phi(m)-1} \equiv 1(\bmod m)$, and $a^{\phi(m)-1}$ is the inverse of a modulo $m$, as desired.

Remark: Corollary 2.19 is useful in solving linear congruences in one variable of the form $a x \equiv b(\bmod m)$ where $(a, m)=1$ i.e. multiply both sides by $a^{\phi(m)-1}$ to solve. Arithmetic Functions and muttiplicativity\\
Definition: An arithmetic function is a function\\
whose domain is the set of positive integers.\\
Example: - Enler $\phi$-function $\phi(n)$ is an orithmetic function.

\begin{itemize}
  \item Let $V(n)$ denote the number of positive divisors of $n$. This is an arithmetic function.
  \item Let $o(n)$ denote the sum of positive divisors of $n$. This is an arithmetic function.
\end{itemize}

Definition : An arithmetic function $f$ is said to be multiplicative if $f(m n)=f(m) f(n)$ whenever $m$ and $n$ are relatively prime positive integers. An arithmetic function $f$ is said to be completely multiplicative if $f(m n)=f(m) f(n)$ for all positive integers $m$ and $n$.

\begin{itemize}
  \item Let $f$ be a multiplicative arithmedic function. If $n$ is a positive integer, then
\end{itemize}


\begin{equation*}
n=p_{1}^{a_{1}} p_{2}^{a_{2}} \cdots p_{r}^{a_{r}} \tag{10}
\end{equation*}


with $p_{1}, \ldots, p_{r}$ distinct prime numbers and $a_{1}, \ldots, a_{r}$ non-negative integers. Since $\left(p_{i}^{a_{i}}, p_{j}^{a_{j}}\right)=1$ if $i \neq j$, we have\\
(米)\\
$f(n)=f\left(p_{1}^{a_{1}} p_{2}^{a_{2}} \cdots p_{r}^{a_{r}}\right)=f\left(p_{1}^{a_{1}}\right) f\left(p_{2}^{a_{2}}\right) \cdots f\left(p_{r}^{a_{n}}\right)$.\\
Multiplicative orithmetic functions are determined by its values at prime powers?\\
Question: What happens for completely multiplicative arithmetic functions $f$ ?\\
Then (*) becomes\\
$f(n)=f\left(p_{1}\right)^{a_{1}} f\left(p_{2}\right)^{a_{2}} \cdots f\left(p_{r}\right)^{a_{n}}$.\\
Completely multiplicative arithmetic functions are determined by its values at primes!


\end{document}